% Introduction — revised 2026-08-09 after the sentence-level fact-check against the data.
% Every quantitative claim is traceable to a measured contrast on the ucr2p_10 cohort
% (ucr_split_w2p, W = 2 x period, per-window z-normalisation, seed 0, cluster as the unit of
% analysis, n = 8 discriminating series of 10). Claims that are bibliographic rather than
% measured are marked TODO-cite and must be substantiated before submission.

\section{Introduction}

Anomaly detection in time series frequently involves data that is distributed by nature:
sensors, machines, or organizations each observe their own series and, for reasons of
confidentiality or ownership, cannot share raw measurements. Each participant is then left
with a fraction of the data a pooled model would see, and collaboration is valuable for that
reason alone: in our experiments a model trained on the pooled shards outperforms the average
shard-local model on all eight discriminating series of our benchmark, and pooling is
precisely what is prohibited. This scarcity is the regime we study. Each client holds a
disjoint slice of the same underlying process, so clients differ in \emph{how much} they have
observed rather than in \emph{what} they observe; cross-client distribution shift is a
distinct problem and we leave it out of scope. Federated learning (FL) offers the standard
alternative; to our knowledge, however, federated work on time-series anomaly detection has
addressed monolithic detectors or continuous-latent generative models.% TODO-cite
How collaboration should be organized in a discrete generative detector -- a pipeline in
which a learned tokenizer feeds a prior over discrete codes -- has not, to our knowledge,
been studied.

We study this question on TimeVQVAE-AD~\citep{lee2024timevqvaead}, a two-stage detector
reporting state-of-the-art results on the UCR benchmark: a VQ-VAE tokenizer maps
time-frequency windows to discrete tokens, and a masked (MaskGIT-style) prior scores
anomalies by the negative log-likelihood it assigns to those tokens. This structure makes
federation non-trivial in a specific way: the prior is a model of the tokenizer's output, so
any federated treatment of the prior presupposes that clients speak a common discrete
vocabulary. We find empirically that they do not. Because each client initialises its own
codebook, code indices are not comparable across clients by construction, and raw agreement
sits at chance ($0.014$ against a chance level of $0.016$); the informative measure is
agreement under the best bijection between indices, and there independently trained
tokenizers reach only $0.34$ -- far above chance, so the partitions do share structure, but
far from the identity, leaving a residual disagreement that no relabeling of the codebook
can close. The consequence is not that sharing the prior is harmful, but that without a
common vocabulary its effect is \emph{unpredictable}: on some series it helps detection and
on others it hurts it, by margins that are far from negligible, and on balance nothing is
gained.

Our central finding is that the obstacle is representation alignment, and that weight
averaging is what removes it. Federating the encoder aligns client tokenizers essentially
exactly, at a token resolution roughly twice that of the other federated variants whose
encoders are trained independently ($1.7$--$2.8\times$ more distinct codes per window),
though still about $30\%$ below the local and centralized tokenizers; and the alignment
comes from averaging the weights rather than from a shared initialization: clients that
start from a common initialization but are never averaged remain as misaligned as
independent ones. Once a common token space exists, sharing the prior changes from
unpredictable to consistently beneficial, and the two ingredients only pay off together: the
same modification of the prior has no net effect on an unaligned backbone, so what alignment
buys is not a larger effect but a directional one -- an interaction rather than a main
effect of either stage. The resulting configuration -- federated tokenizer, shared prior
body with lightweight per-client heads, no exchange of raw series -- is, to our knowledge,
the first federated adaptation of a discrete tokenizer--prior detector that improves over
local models trained to convergence under the same stopping rule. The effect is concentrated
rather than uniform: on three series the gain exceeds the spread we measure across training
protocols ($+0.30$ to $+0.40$ AUPRC), on one series the \emph{loss} does too ($-0.43$), and
on the remaining four the two are indistinguishable. It does not close the gap to centralized
training; we therefore decompose the residual gap, attributing it between the tokenizer and
the prior stage, and characterize the failure cases, including the series on which the
federated detector does not recover the local baseline.

All comparisons are paired per series with the client cluster as the unit of analysis, since
the clients of a cluster share a test set and their scores are strongly correlated. The
federated arms stop on a validation-driven convergence criterion, every cell halting before
its round ceiling, whereas the local and centralized baselines stop on a hard step budget in
about half the cells of the main configuration; we therefore re-run both baselines under a
raised budget and report both. Results are read against noise thresholds estimated from
replicate runs and against the sensitivity of each baseline to its own training protocol,
and are calibrated against a zero-parameter baseline that several federated variants fail to
beat.

In summary, our contributions are:

\begin{enumerate}
    \item an empirical decomposition of where to federate a two-stage discrete generative
    anomaly detector, showing that tokenizer alignment is a necessary precondition and that
    the collaboration gain arises from sharing the prior once that precondition holds --
    without it, the same sharing moves detection as much but in inconsistent directions, for
    no net gain;
    \item a mechanism analysis showing that encoder weight averaging, and not a shared
    initialization, is what produces the aligned vocabulary that naive federation lacks, at
    roughly twice the token resolution of the federated variants that leave encoders local;
    \item a federated configuration that improves over converged local training on the UCR
    anomaly benchmark without sharing raw time series, with the magnitude and the extent of
    the improvement reported per series against measured variability; and
    \item a diagnostic that separates the remaining tokenizer-side and prior-side gaps to
    centralized training, indicating where future federated designs should invest.
\end{enumerate}
